\textbf{1. Kiến thức cần nhớ}
\begin{itemize}
\item \textbf{Căn bậc ba} của một số thực $a$ là số thực $x$ sao cho $x^3 = a$, kí hiệu là $\sqrt[3]{a}$.
\begin{itemize}
\item Mỗi số thực $a$ đều có \textbf{duy nhất} một căn bậc ba.
\item Căn bậc ba của số dương là số dương: $\sqrt[3]{a} > 0$ khi $a > 0$.
\item Căn bậc ba của số âm là số âm: $\sqrt[3]{a} < 0$ khi $a < 0$.
\item $\sqrt[3]{0} = 0$.
\item Căn thức bậc ba $\sqrt[3]{A}$ \textbf{luôn xác định với mọi biểu thức $A$} (không cần điều kiện $A \ge 0$).
\end{itemize}
\item \textbf{Tính chất}:
$$\boxed{(\sqrt[3]{a})^3 = \sqrt[3]{a^3} = a}$$
$$\boxed{a < b \iff \sqrt[3]{a} < \sqrt[3]{b}}$$
$$\boxed{\sqrt[3]{ab} = \sqrt[3]{a} \cdot \sqrt[3]{b}};\quad \boxed{\sqrt[3]{\dfrac{a}{b}} = \dfrac{\sqrt[3]{a}}{\sqrt[3]{b}}\quad (b \ne 0)}$$
\end{itemize}

\textbf{2. Phương pháp giải các dạng}

%% DANG 9D3?10-1
\textbf{Dạng 1. Tính căn bậc ba, so sánh}
\begin{enumerate}
\item Nhận diện các lập phương hoàn chỉnh: $1^3 = 1, 2^3 = 8, 3^3 = 27, 4^3 = 64, 5^3 = 125\dots$
\item Đưa biểu thức dưới dấu căn về dạng $A^3$ rồi dùng hằng đẳng thức $\sqrt[3]{A^3} = A$.
\item So sánh: $a < b \iff \sqrt[3]{a} < \sqrt[3]{b}$ (hoặc lập phương hai vế để so sánh).
\end{enumerate}

%% DANG 9D3?10-2
\textbf{Dạng 2. Biến đổi, rút gọn biểu thức chứa căn bậc ba}
\begin{enumerate}
\item Sử dụng các hằng đẳng thức lập phương:
$$(a \pm b)^3 = a^3 \pm 3a^2b + 3ab^2 \pm b^3;\quad a^3 \pm b^3 = (a \pm b)(a^2 \mp ab + b^2).$$
\item Nhân lượng liên hợp bậc ba để trục căn thức ở mẫu:
$$\dfrac{1}{\sqrt[3]{a} \pm \sqrt[3]{b}} = \dfrac{\sqrt[3]{a^2} \mp \sqrt[3]{ab} + \sqrt[3]{b^2}}{a \pm b}.$$
\end{enumerate}

%% DANG 9D3?10-3
\textbf{Dạng 3. Phương trình chứa căn bậc ba đơn giản}
\begin{enumerate}
\item Phương trình dạng $\sqrt[3]{f(x)} = a \iff \boxed{f(x) = a^3}$ (không cần đặt điều kiện xác định).
\item Phương trình dạng $\sqrt[3]{f(x)} = \sqrt[3]{g(x)} \iff \boxed{f(x) = g(x)}$.
\item Lập phương hai vế và giải phương trình đại số thu được.
\end{enumerate}
